{\small\textbf{Monetary and macroprudential policy (Zheng, 1:05) [18:55–20:00]}}\par\smallskip
Thank you. The first wider implication is the interaction with monetary policy.
\begin{itemize}\setlength\itemsep{0pt}
\item The worry was a \textbf{double squeeze}: the ECB raising rates while DNB raised capital requirements.
\item As we just saw, Dutch lending rates did not diverge from their peers, and bank profits rose with rates.
\end{itemize}
\par\smallskip
The ECB argues that building buffers early "helps monetary policy focus on its primary objective of price stability". \textcolor{navy}{\textbf{[def]}} Macroprudential policy targets the stability of the financial system as a whole, and it does not have to fight inflation.
\par\smallskip
For a euro-area country this matters even more. There is one monetary policy for twenty countries, but macroprudential policy is national. The countercyclical buffer is the Dutch-specific lever, and it is the insurance in case a rate shock goes wrong.
